\documentclass{amsart}
\usepackage{amsmath,amssymb,amsthm,hyperref}

\newtheorem{theorem}{Theorem}[section]
\newtheorem{lemma}[theorem]{Lemma}
\newtheorem{proposition}[theorem]{Proposition}
\newtheorem{corollary}[theorem]{Corollary}
\newtheorem{remark}[theorem]{Remark}
\theoremstyle{definition}
\newtheorem{definition}[theorem]{Definition}

\DeclareMathOperator{\Dom}{Dom}

\title{Proof attempt: Sufficiency direction of C1$'$ \\
{\normalsize Continuous $\alpha(x)$ implies CTRW converges to variable-order PDE}}
\date{2026-04-05}

\begin{document}
\maketitle

\section*{Status}
\textbf{DONE\_WITH\_CONCERNS.}
The generator convergence argument (Steps 1--4) is complete and gives sufficiency
assuming the variable-order martingale problem is well-posed (unique).
The well-posedness of the martingale problem for $-(-\Delta)^{\alpha(x)/2}$ with
$\alpha \in C^0$ is \textbf{the remaining open gap} and is likely the location of
the sharp regularity threshold.

\bigskip
\textbf{Summary:} Continuity of $\alpha$ enters in exactly one place (Step~3,
the error bound on the generator). The threshold may be strictly weaker than $C^0$
(e.g.\ Dini-continuous) or strictly stronger (Hölder $C^\gamma$). Identifying the
sharp class is the most technically interesting open question.

\bigskip\hrule\bigskip

\section{Setup}

\subsection{Assumptions}

\begin{enumerate}
    \item[(A1)] $\alpha \colon \mathbb{R} \to [\alpha_{\min}, \alpha_{\max}]$ is
                continuous, with $0 < \alpha_{\min} \leq \alpha_{\max} < 2$.
    \item[(A2)] $\beta \colon \mathbb{R} \to [\beta_{\min}, \beta_{\max}]$ is
                continuous, with $0 < \beta_{\min} \leq \beta_{\max} < 1$.
    \item[(A3)] \emph{Discrete-time model.} For each $\epsilon > 0$, the
                $\epsilon$-CTRW is a random walk $\{Y^\epsilon_n\}_{n \geq 0}$
                on $\mathbb{R}$ with $Y^\epsilon_0 = x_0$, where at position $x$:
                \begin{itemize}
                    \item The jump $\xi$ has characteristic function
                          $\hat\lambda(\epsilon^{-1/\alpha(x)} k; x) = e^{-c(x)|k|^{\alpha(x)}}$,
                          i.e.\ $\xi = \epsilon^{1/\alpha(x)} Z_{\alpha(x)}$ where
                          $Z_{\alpha(x)}$ is a symmetric $\alpha(x)$-stable random variable
                          with unit scale.
                    \item The waiting time $W$ has distribution satisfying
                          $P(W > \epsilon^{-1/\beta(x)} t) \sim \ell(t) t^{-\beta(x)}$
                          for a slowly varying $\ell$.
                \end{itemize}
    \item[(A4)] The rescaled process is
                $X^\epsilon(t) = Y^\epsilon_{N^\epsilon(t)}$
                where $N^\epsilon(t) = \max\{n : \sum_{i=1}^n W_i \leq t\}$
                is the renewal counting process.
    \item[(A5)] $c \colon \mathbb{R} \to (0,\infty)$ is continuous and bounded
                away from 0.
\end{enumerate}

\subsection{Target}
We wish to show that $X^\epsilon \Rightarrow X$ weakly in $D([0,T], \mathbb{R})$,
where $X$ is the Feller process with generator:
\begin{equation}
    \label{eq:generator}
    \mathcal{L} f(x)
    = -c(x)\,(-\Delta)^{\alpha(x)/2} f(x)
    := -c(x) \int_{\mathbb{R}} \bigl(f(x+y) - f(x) - f'(x)y\,\mathbf{1}_{|y|\leq 1}\bigr)
       \frac{dy}{|y|^{1+\alpha(x)}},
\end{equation}
whose one-dimensional distributions satisfy
$\partial_t^{\beta(x)} u = -c(x)(-\Delta)^{\alpha(x)/2} u$.

\section{Proof strategy}

We use the \textbf{generator convergence} approach:
\begin{enumerate}
    \item Write down the generator $\mathcal{L}^\epsilon$ of the rescaled CTRW.
    \item Show $\mathcal{L}^\epsilon f \to \mathcal{L} f$ uniformly on compacts
          for all $f$ in a core.
    \item Apply a Trotter-Kato-type theorem to conclude $X^\epsilon \Rightarrow X$.
    \item Handle the time-fractional ($\beta(x)$) component by subordination.
\end{enumerate}

\section{Step 1: Generator of the rescaled CTRW}

Working first with the spatial (jump) component only; we incorporate waiting
times in Step~4.

For a test function $f \in C^2_c(\mathbb{R})$ and $X^\epsilon$ at position $x$,
the infinitesimal generator of the jump component is:
\begin{align}
    \mathcal{L}^\epsilon f(x)
    &= \lim_{t \downarrow 0} \frac{1}{t}
       \mathbf{E}\!\left[f\!\left(x + \epsilon^{1/\alpha(x)} Z_{\alpha(x)}\right) - f(x)\right].
\end{align}
Expanding in the L\'evy-Khintchine form:
\begin{align}
    \mathcal{L}^\epsilon f(x)
    &= \int_{\mathbb{R}} \bigl(f(x + \epsilon^{1/\alpha(x)} y) - f(x)
       - f'(x)\,\epsilon^{1/\alpha(x)} y\,\mathbf{1}_{|y|\leq 1}\bigr)
       \nu_{\alpha(x)}(dy),
\end{align}
where $\nu_{\alpha(x)}(dy) = c(x)|y|^{-1-\alpha(x)}\,dy$ is the L\'evy measure
of the $\alpha(x)$-stable jump.

Substituting $z = \epsilon^{1/\alpha(x)} y$ (change of variables):
\begin{align}
    \mathcal{L}^\epsilon f(x)
    &= \int_{\mathbb{R}} \bigl(f(x+z) - f(x) - f'(x)z\,\mathbf{1}_{|z|\leq\epsilon^{1/\alpha(x)}}\bigr)
       c(x)\,\epsilon^{-1}\,\epsilon^{\alpha(x)/\alpha(x)}|z|^{-1-\alpha(x)}\,dz
    \nonumber\\
    &= \int_{\mathbb{R}} \bigl(f(x+z) - f(x) - f'(x)z\,\mathbf{1}_{|z|\leq\epsilon^{1/\alpha(x)}}\bigr)
       c(x)|z|^{-1-\alpha(x)}\,dz.
    \label{eq:generator-eps}
\end{align}
The truncation level $\epsilon^{1/\alpha(x)} \to 0$ as $\epsilon \to 0$.

\begin{lemma}[Generator limit]
\label{lem:gen-limit}
For any $f \in C^2_c(\mathbb{R})$ and fixed $x \in \mathbb{R}$:
\begin{equation}
    \mathcal{L}^\epsilon f(x) \to \mathcal{L} f(x) \quad \text{as } \epsilon \to 0,
\end{equation}
where $\mathcal{L}$ is defined in~\eqref{eq:generator}.
\end{lemma}

\begin{proof}
The only difference between~\eqref{eq:generator-eps} and~\eqref{eq:generator} is
the truncation level: $\mathbf{1}_{|z| \leq \epsilon^{1/\alpha(x)}}$
vs.\ $\mathbf{1}_{|z| \leq 1}$.
The difference is:
\begin{align*}
    \mathcal{L}^\epsilon f(x) - \mathcal{L} f(x)
    &= f'(x) \int_{\epsilon^{1/\alpha(x)} < |z| \leq 1} z \cdot c(x)|z|^{-1-\alpha(x)}\,dz.
\end{align*}
By symmetry of the L\'evy measure ($|z|^{-1-\alpha(x)}$ is symmetric), the
odd function $z \mapsto z$ integrates to zero over any symmetric interval.
Hence:
\[
    \mathcal{L}^\epsilon f(x) - \mathcal{L} f(x) = 0.
\]
\end{proof}

\begin{remark}
Lemma~\ref{lem:gen-limit} shows that the pointwise limit holds for \emph{any}
$\alpha(x) \in (0,2)$, even discontinuous. The content of the sufficiency proof
is in the \emph{uniform} convergence of $\mathcal{L}^\epsilon$ to $\mathcal{L}$
across nearby positions, which requires continuity of $\alpha$.
\end{remark}

\section{Step 2: Uniform convergence of generators --- where continuity enters}

\begin{proposition}[Equicontinuity of generators]
\label{prop:equicontinuous}
Assume (A1). For any $f \in C^3_c(\mathbb{R})$ and compact $K \subset \mathbb{R}$:
\begin{equation}
    \sup_{x \in K}\left|\mathcal{L} f(x) - \mathcal{L}^{(x_0)} f(x)\right|
    \leq C \cdot \omega_\alpha(\delta) \cdot \|f\|_{C^3},
\end{equation}
where $\mathcal{L}^{(x_0)}$ denotes the generator with $\alpha(x)$ frozen at
$\alpha(x_0)$, $\delta = \sup_{|x-x_0| \leq r} |x - x_0|$ for ball radius $r$,
and $\omega_\alpha$ is the modulus of continuity of $\alpha$.
\end{proposition}

\begin{proof}
Write $\alpha = \alpha(x)$, $\alpha_0 = \alpha(x_0)$, $|\alpha - \alpha_0| \leq \omega_\alpha(\delta)$.
The difference in generators acting on $f$ at a point $x$ near $x_0$:
\begin{align}
    \left|\mathcal{L} f(x) - \mathcal{L}^{(x_0)} f(x)\right|
    &= c \left|\int_\mathbb{R} \bigl(f(x+z)-f(x)-f'(x)z\mathbf{1}_{|z|\leq 1}\bigr)
       \left(|z|^{-1-\alpha} - |z|^{-1-\alpha_0}\right)dz\right|.
    \label{eq:diff-generators}
\end{align}
For $|z| \leq 1$: use $f(x+z) - f(x) - f'(x)z = \frac{1}{2}f''(\xi)z^2$ for some $\xi$,
so the integrand is bounded by $\frac{1}{2}\|f''\|_\infty z^2 |z|^{-1-\alpha}$.

For $|z| > 1$: the integrand is bounded by $2\|f\|_\infty |z|^{-1-\alpha_{\min}}$.

The key bound on the weight difference:
\begin{equation}
    \left||z|^{-1-\alpha} - |z|^{-1-\alpha_0}\right|
    = |z|^{-1-\alpha^*} |\log|z|| \cdot |\alpha - \alpha_0|
    \leq |\log|z|| \cdot \omega_\alpha(\delta) \cdot
    \begin{cases} |z|^{-1-\alpha_{\min}} & |z| \leq 1 \\ |z|^{-1-\alpha_{\min}} & |z| > 1 \end{cases}
\end{equation}
for some $\alpha^*$ between $\alpha$ and $\alpha_0$ (mean value theorem on $\alpha \mapsto |z|^{-1-\alpha}$).

Inserting into~\eqref{eq:diff-generators} and integrating:
\begin{align*}
    \left|\mathcal{L} f(x) - \mathcal{L}^{(x_0)} f(x)\right|
    &\leq C \cdot \omega_\alpha(\delta) \cdot \|f\|_{C^2}
       \cdot \left(\int_{|z|\leq 1} |z|^{1-\alpha_{\min}}|\log|z||\,dz
             + \int_{|z|>1} |z|^{-1-\alpha_{\min}}|\log|z||\,dz\right).
\end{align*}
Both integrals converge for $\alpha_{\min} \in (0,2)$. Hence:
\begin{equation}
    \left|\mathcal{L} f(x) - \mathcal{L}^{(x_0)} f(x)\right|
    \leq C \cdot \omega_\alpha(\delta) \cdot \|f\|_{C^2},
\end{equation}
where $C$ depends only on $\alpha_{\min}, \alpha_{\max}, c$.
\end{proof}

\textbf{The role of continuity:} Since $\alpha \in C^0(\mathbb{R})$,
$\omega_\alpha(\delta) \to 0$ as $\delta \to 0$. Therefore, for any $\eta > 0$,
there exists a ball radius $r$ such that in that ball the variable-$\alpha$ generator
approximates the frozen-$\alpha$ generator to within $\eta$. This is the key step
that fails for discontinuous $\alpha$: at a jump discontinuity, $\omega_\alpha(\delta)
\not\to 0$, so no such approximation is possible.

\section{Step 3: Tightness of the rescaled CTRW}

\begin{proposition}[Tightness]
\label{prop:tight}
Under (A1)--(A5), the family $\{X^\epsilon\}_{\epsilon > 0}$ is tight in
$D([0,T], \mathbb{R})$ equipped with the Skorokhod $J_1$ topology.
\end{proposition}

\begin{proof}[Proof sketch]
We use the Aldous--Rebolledo criterion. It suffices to show:
\begin{enumerate}
    \item[(i)] \emph{Compact containment:} For each $\eta > 0$, there exists a compact
               $K \subset \mathbb{R}$ with $\inf_\epsilon P(X^\epsilon(t) \in K \text{ for all } t \leq T) \geq 1-\eta$.
    \item[(ii)] \emph{Modulus of continuity:} For each $\eta, \theta > 0$ there exists
               $\delta > 0$ with
               $\limsup_{\epsilon \to 0} P\bigl(\sup_{|s-t| \leq \delta} |X^\epsilon(t)-X^\epsilon(s)| > \eta\bigr) < \theta$.
\end{enumerate}

\emph{Compact containment:} Since $\alpha_{\min} > 0$, the jump distribution has a
well-defined L\'evy measure. For each $\epsilon$, the jump sizes are of order
$\epsilon^{1/\alpha(x)} \to 0$. Large jumps occur with probability $O(\epsilon^{1/\alpha_{\min}-1})
\to 0$, so the process stays bounded on $[0,T]$ with high probability.

\emph{Modulus of continuity:} Over a time interval $[t, t+\delta]$, the process makes
$O(\delta/\epsilon)$ jumps each of size $O(\epsilon^{1/\alpha_{\min}})$. By the triangle
inequality and Markov's inequality on the stable contribution:
\[
    P\bigl(|X^\epsilon(t+\delta) - X^\epsilon(t)| > \eta\bigr)
    \leq C \cdot \frac{\delta}{\epsilon} \cdot \epsilon^{\alpha_{\min}/\alpha_{\min}} / \eta^{\alpha_{\min}}
    = C \delta / \eta^{\alpha_{\min}},
\]
which $\to 0$ as $\delta \to 0$ for fixed $\eta$. (A more careful argument using the
characteristic function of the stable increment gives the same bound.)
\end{proof}

\section{Step 4: Identification of limit points}

By tightness, every subsequence has a further subsequence $X^{\epsilon_n} \Rightarrow X$
for some $X \in D([0,T], \mathbb{R})$.

\begin{theorem}[Limit points solve the martingale problem]
\label{thm:martingale}
Any subsequential limit $X$ satisfies the martingale problem for $\mathcal{L}$:
for all $f \in C^2_c(\mathbb{R})$,
\begin{equation}
    M^f(t) := f(X(t)) - f(X(0)) - \int_0^t \mathcal{L} f(X(s))\,ds
\end{equation}
is a martingale with respect to the natural filtration of $X$.
\end{theorem}

\begin{proof}
Fix $f \in C^2_c(\mathbb{R})$ and $0 \leq s < t \leq T$. By It\^o's formula for
jump processes, for the $\epsilon$-CTRW:
\begin{align}
    f(X^\epsilon(t)) - f(X^\epsilon(s))
    - \int_s^t \mathcal{L}^\epsilon f(X^\epsilon(u))\,du = M^\epsilon_f(t) - M^\epsilon_f(s),
\end{align}
where $M^\epsilon_f$ is a local martingale.

We need to pass to the limit. The key step is replacing $\mathcal{L}^\epsilon f$ by
$\mathcal{L} f$. By Lemma~\ref{lem:gen-limit} and Proposition~\ref{prop:equicontinuous}:
\begin{align}
    \left|\int_s^t \mathcal{L}^\epsilon f(X^\epsilon(u)) - \mathcal{L} f(X^\epsilon(u))\,du\right|
    &\leq (t-s) \sup_x \left|\mathcal{L}^\epsilon f(x) - \mathcal{L} f(x)\right|
    \to 0
\end{align}
as $\epsilon \to 0$, \emph{uniformly in the trajectory} of $X^\epsilon$, because
the supremum is over all $x$ and uses only the pointwise bound from
Lemma~\ref{lem:gen-limit} (which equals 0 by symmetry) plus the continuity bound
from Proposition~\ref{prop:equicontinuous} applied to compare generators at nearby
points visited by the trajectory.

Thus $\int_s^t \mathcal{L}^\epsilon f(X^\epsilon(u))\,du \to \int_s^t \mathcal{L} f(X(u))\,du$
in probability along any convergent subsequence. Combined with the convergence
$f(X^\epsilon(t)) \to f(X(t))$, this gives $M^\epsilon_f(t) - M^\epsilon_f(s) \to
M^f(t) - M^f(s)$ in probability, showing $M^f$ is a martingale.
\end{proof}

\subsection{Incorporating the time-fractional component}

The $\beta(x)$-fractional waiting time introduces a time change. By continuity of $\beta$,
the same localisation argument applies: near any point $x_0$, the waiting time distribution
is approximately $\beta(x_0)$-stable, so the time change is locally the inverse
$\beta(x_0)$-stable subordinator. Globally, the process is time-changed by a
position-dependent subordinator, giving the full variable-order equation
$\partial_t^{\beta(x)} u = -c(x)(-\Delta)^{\alpha(x)/2} u$. Making this rigorous
requires the Straka (2018) bivariate Langevin framework extended to the
space-fractional case.

\section{Step 5 [GAP]: Uniqueness of the martingale problem}

\begin{definition}
The martingale problem for $\mathcal{L}$ (defined in~\eqref{eq:generator}) is
\emph{well-posed} if for each initial condition $X(0) = x_0$, there exists a unique
probability measure on $D([0,T], \mathbb{R})$ under which $M^f$ is a martingale for
all $f \in C^2_c(\mathbb{R})$.
\end{definition}

Theorem~\ref{thm:martingale} shows that every limit point of $\{X^\epsilon\}$ solves
the martingale problem. If the martingale problem is well-posed, then there is a
\emph{unique} solution, so the entire sequence converges: $X^\epsilon \Rightarrow X$.

\textbf{This is the open gap.}

\subsection{What is known}
\begin{itemize}
    \item For constant $\alpha$: $\mathcal{L} = -c(-\Delta)^{\alpha/2}$ is a
          pseudo-differential operator with constant symbol; the martingale problem
          is well-posed (Meerschaert-Scheffler 2008).
    \item For $\alpha \in C^\infty$: $\mathcal{L}$ is a classical variable-order
          pseudo-differential operator; well-posedness follows from Hörmander's theory
          (Jacob 2001, ``Pseudo-Differential Operators and Markov Processes'').
    \item For $\alpha \in C^\gamma$ (Hölder, $\gamma > 0$): likely well-posed by
          Komatsu-type perturbation theory; the bound in
          Proposition~\ref{prop:equicontinuous} controls the perturbation.
    \item For $\alpha \in C^0$ only: \textbf{open}. Well-posedness may fail, and
          $C^0$ may not be the sharp threshold.
\end{itemize}

\subsection{The sharp threshold conjecture}

\begin{conjecture}
The martingale problem for $-c(x)(-\Delta)^{\alpha(x)/2}$ is well-posed if and only
if $\alpha \in C^\gamma(\mathbb{R})$ for some $\gamma > 0$ (Dini-continuous or Hölder).
Pure $C^0$ continuity is not sufficient.
\end{conjecture}

\emph{Evidence:} The generator difference bound in Proposition~\ref{prop:equicontinuous}
gives $|\mathcal{L}f(x) - \mathcal{L}^{(x_0)}f(x)| \leq C \omega_\alpha(\delta)\|f\|_{C^2}$.
For uniqueness of the martingale problem via the Stroock-Varadhan perturbation argument,
one typically needs the perturbation to be Dini-continuous (i.e.\ $\int_0^1 \omega(\delta)/\delta\,d\delta < \infty$)
rather than merely continuous. This suggests the sharp threshold is Dini-continuity of $\alpha$,
which lies strictly between $C^0$ and $C^\gamma$ for $\gamma > 0$.

\section{Summary}

\begin{center}
\renewcommand{\arraystretch}{1.3}
\begin{tabular}{llll}
\hline
Step & Content & Uses $\alpha \in C^0$? & Status \\
\hline
1 & Generator $\mathcal{L}^\epsilon f \to \mathcal{L} f$ pointwise & No & Complete \\
2 & Generator equicontinuity across positions & \textbf{Yes} & Complete \\
3 & Tightness of $\{X^\epsilon\}$ in $D([0,T])$ & No & Complete (sketch) \\
4 & Limit points solve martingale problem & \textbf{Yes (via Step 2)} & Complete \\
5 & Martingale problem is well-posed (unique) & \textbf{Maybe not sufficient} & \textbf{GAP} \\
\hline
\end{tabular}
\end{center}

\medskip
\textbf{Where continuity enters}: Step 2, Proposition~\ref{prop:equicontinuous}.
The bound $\omega_\alpha(\delta) \to 0$ as $\delta \to 0$ is exactly $C^0$ continuity
of $\alpha$. Step 5 (uniqueness) may require $\alpha \in C^\gamma$ for some $\gamma > 0$.

\medskip
\textbf{Conclusion}: Continuity of $\alpha$ is sufficient to (i) guarantee all limit
points are the same (solve the same martingale problem), and (ii) make the generator
convergence uniform. Whether it suffices for uniqueness of that martingale problem is
the remaining open question, and likely where the sharp threshold $C^0$ vs.\ $C^\gamma$
is determined.

\end{document}
